\PID Период осциловања тела масе $m$ обешеног о опругу коефицијента еластичности $k$ 
дат је у облику 
$T = 2\uppi\sqrt{\dfrac{m}{k}}$. Ако су познате релативне грешке  
$\updelta m = 5\%$ и $\updelta k = 5\%$, израчунати релативну грешку 
периода $\updelta T$. 

\RESENJE 

Релативна грешка одређује се као релативни удео апсолутне грешке у односу на вредност мерене величине, односно 
$\updelta T = \dfrac{\Delta T}{T}$. Апсолутна грешка може се одредити формалним поступком који је дат изразом
\eqref{eq:def_aps_greska}, али је погодније применити исту технику која је илустрована у задатку 
\ref{z:proizvod}. Уопштење правила у производу се може добити применом својства логаритма
$\ln x^a = a\ln x$, $a = \rm const$, водећи рачуна о најгорем случају, па је 
\begin{equation}
    \updelta(x_1^{a_1}x_2^{a_2}\cdots x_n^{a_n})
    = |a_1| \, \updelta x_1 + \cdots |a_n| \, \updelta x_n,
\end{equation}
па се применом а израз за период има $\updelta T = \dfrac{1}{2} \updelta m + \dfrac{1}{2} \updelta k = 5\%$.